% 附录C 玻色与费米相干态
\chapter{Bose 与 Fermi 相干态}

相干态构成超完备基底，以连续参数或变量标记。相干态基底中的单位分解常用于计算算符的矩阵元与迹。这里定义这些态，并给出直接由定义得到的若干有用恒等式。

\section{复积分}

Bose 相干态由复矢量参数化：

\begin{equation*}
\mathbf {z} = \left( z _ {1}, z _ {2}, \dots , z _ {\mathcal {N}} \right) \, , \quad z _ {i} = x _ {i} + iy _ {i} ,
\tag{C.1}
\end{equation*}

$x_{i}, y_{i}$ 为实数，$\mathcal{N}$ 是单粒子态数目。复积分定义为

\begin{equation*}
\int d ^ {2} \mathbf {z} \equiv \int_{- \infty} ^ {\infty} \prod_ {i} \left( \frac {dx _ {i} \, dy _ {i}}{\pi} \right).
\tag{C.2}
\end{equation*}

对单变量 $z$ 的复积分有一个有用恒等式：

\begin{equation*}
\frac {1}{n !} \int d ^ {2} \mathbf {z} \left( z ^ {*} \right) ^ {n} z ^ {m} \exp \left( - z ^ {*} z \right) = \delta_{n, m}.
\tag{C.3}
\end{equation*}

利用这一定义可知：对任何厄米部分只有正本征值的复矩阵 $G$，Gauss 积分为

\begin{equation*}
\int d ^ {2} \mathbf {z} \, e ^ {- \mathbf {Z} ^ {*} G \mathbf {Z} - \mathbf {Z} _ {a} ^ {*} \mathbf {Z} - \mathbf {Z} ^ {*} \mathbf {Z} _ {b}} = \det | G | ^ {- 1} e ^ {\mathbf {Z} _ {a} ^ {*} G ^ {- 1} \mathbf {Z} _ {b}} ,
\tag{C.4}
\end{equation*}

$\mathbf{z}_a^*, \mathbf{z}_b$ 是任意复矢量。

\section{Grassmann 变量}

Fermi 相干态由共轭的 Grassmann 矢量对参数化：

\begin{equation*}
\mathbf {z} = \left( z _ {1}, z _ {2}, \dots , z _ {\mathcal {N}} \right) \, , \quad \mathbf {z} ^ {*} = \left( z _ {1} ^ {*}, z _ {2} ^ {*}, \dots , z _ {\mathcal {N}} ^ {*} \right) ,
\tag{C.5}
\end{equation*}

$\mathcal{N}$ 是单粒子态数目。这些变量是可以相加、相乘的算符：它们与任何标量 $c$ 对易，但彼此反对易：

\begin{align*}
cz & = zc ,\\
z _ {i} z _ {j} & = - z _ {j} z _ {i} ,\\
z _ {i} z _ {j} ^ {*} & = - z _ {j} ^ {*} z _ {i} .
\tag{C.6}
\end{align*}

（Grassmann 变量的共轭视为另一个变量。）(C.6) 蕴含

\begin{equation*}
z _ {i} ^ {2} = 0.
\tag{C.7}
\end{equation*}

因此任何 Grassmann 变量的函数都可以展开为有限和。

“Grassmann 积分”是一种计数操作，作用于 Grassmann 变量的乘积：

\begin{equation*}
\int d z _ {1} \prod_{i = 1} ^ {\mathcal {N}} z _ {i} ^ {n _ {i}} \prod_{j = 1} ^ {\mathcal {N}} \left( z _ {j} ^ {*} \right) ^ {m _ {j}} c = n _ {1} \prod_{i = 2} ^ {\mathcal {N}} z _ {i} ^ {n _ {i}} \prod_{j = 1} ^ {\mathcal {N}} \left( z _ {j} ^ {*} \right) ^ {m _ {j}} c ,
\tag{C.8}
\end{equation*}

$n_{i}, m_{j} = 0, 1$，$c$ 是任意标量。对一般的 Grassmann 乘积，先把要积分的变量反对易到最左端，同时得到置换的整体符号。由于 $z_{1}$ 与 $z_{1}^{*}$ 视为两个不同变量，把对 $d^{2}z_{1}$ 的积分定义为

\begin{equation*}
\int d ^ {2} z _ {1} \, O ( \mathbf {z} ^ {*}, \mathbf {z} ) \equiv \int d z _ {1} ^ {*} \left[ \int d z _ {1} \, O ( \mathbf {z} ^ {*}, \mathbf {z} ) \right] ,
\tag{C.9}
\end{equation*}

$O$ 是乘积之和。由 (C.8)，每个乘积除非恰好含 $z_{1}$ 的一次幂与 $z_{1}^{*}$ 的一次幂，其贡献为零。由上述定义，对一个变量及其复共轭，Grassmann 的 Gauss 积分为

\begin{equation*}
\int d ^ {2} z \, \exp \left( - z ^ {*} z \right) z ^ {m} \left( z ^ {*} \right) ^ {n} = \delta_{mn} \, , \quad n, m = 0, 1 ,
\tag{C.10}
\end{equation*}

形式上与复积分 (C.3) 相似。多维 Gauss 积分为

\begin{equation*}
\int d ^ {2} \mathbf {z} \, \exp \left( - \mathbf {z} ^ {*} G \mathbf {z} + \mathbf {z} _ {a} ^ {*} \mathbf {z} + \mathbf {z} ^ {*} \mathbf {z} _ {b} \right) = \det | G | \, \exp \left( \mathbf {z} _ {a} ^ {*} G ^ {- 1} \mathbf {z} _ {b} \right) ,
\tag{C.11}
\end{equation*}

$z_{a}^{*}, z_{b}$ 是 Grassmann 矢量。(C.4) 与 (C.11) 可统一为单一表达式

\begin{equation*}
\int d ^ {2} \mathbf {z} \, \exp \left( - \mathbf {z} ^ {*} G \mathbf {z} + \mathbf {z} _ {a} ^ {*} \mathbf {z} + \mathbf {z} ^ {*} \mathbf {z} _ {b} \right) = \det | G | ^ {- \eta} \exp \left( \mathbf {z} _ {a} ^ {*} G ^ {- 1} \mathbf {z} _ {b} \right) ,
\tag{C.12}
\end{equation*}

其中

\begin{equation*}
\eta = \left\{ \begin{array}{ll} 1 & \text{玻色子} \\ - 1 & \text{费米子} \end{array} \right. .
\tag{C.13}
\end{equation*}

\section{相干态}

相干态定义为

\begin{equation*}
\left| \mathbf {z} \right\rangle = \exp \left( \mathbf {za} ^ {\dagger} \right) \left| 0 \right\rangle ,
\tag{C.14}
\end{equation*}

$\mathbf{a}$ 是 Bose（Fermi）产生算符的矢量，$\mathbf{z}$ 是复数（Grassmann 变量）的矢量。相干态是湮灭算符的本征态：

\begin{equation*}
a _ {i} \left| \mathbf {z} \right\rangle = z _ {i} \left| \mathbf {z} \right\rangle .
\tag{C.15}
\end{equation*}

两个相干态的交叠为

\begin{equation*}
\left\langle \mathbf {z} \middle| \mathbf {z} ' \right\rangle = e ^ {\mathbf {z} ^ {*} \mathbf {z} '}.
\tag{C.16}
\end{equation*}

于是基底 $\{|z\rangle\}$ 不正交但张成 Fock 空间——由积分给出的单位分解

\begin{equation*}
\int d ^ {2} z \, e ^ {- \mathbf {z} ^ {*} \mathbf {z}} \left| \mathbf {z} \right\rangle \left\langle \mathbf {z} \right| = I ,
\tag{C.17}
\end{equation*}

即可看出，它由 (C.3) 与 (C.10) 得到。

任何能用产生与湮灭算符展开的算符都可以正规排序，即写成每个乘积中产生算符位于湮灭算符左侧的和式。我们把讨论限于形如

\begin{align*}
\mathcal {O} [ \mathbf {a} ^ {\dagger}, \mathbf {a} ] & = \sum_{\{n _ {i}, m _ {i} \}} \mathcal {O}_{\{n _ {i}, m _ {i} \}} \left[ \left( a _ {1} ^ {\dagger} \right) ^ {n _ {1}} \dots \left( a _ {\mathcal {N}} ^ {\dagger} \right) ^ {n _ {\mathcal {N}}} a _ {1} ^ {m _ {1}} \dots a _ {\mathcal {N}} ^ {m _ {\mathcal {N}}} \right]\\
& = : \mathcal {O} : ,
\tag{C.18}
\end{align*}

的正规排序算符，对费米子 $n_{i}, m_{j} \leq 1$。用 (C.15)，任意正规排序算符在相干态之间的矩阵元为

\begin{equation*}
\left\langle \mathbf {z} \left| \mathcal {O} [ \mathbf {a} ^ {\dagger}, \mathbf {a} ] \right| \mathbf {z} ' \right\rangle = O ( \mathbf {z} ^ {*}, \mathbf {z} ' ) \, e ^ {\mathbf {z} ^ {*} \mathbf {z} '} ,
\tag{C.19}
\end{equation*}

其中

\begin{equation*}
O \left( \mathbf {z} ^ {*}, \mathbf {z} ' \right) = \sum_{\{n _ {i}, m _ {i} \}} O_{\{n _ {i}, m _ {i} \}} \left[ \left( z _ {1} ^ {*} \right) ^ {n _ {1}} \dots \left( z _ {\mathcal {N}} ^ {*} \right) ^ {n _ {\mathcal {N}}} z _ {1} ^ {m _ {1}} \dots z _ {\mathcal {N}} ^ {m _ {\mathcal {N}}} \right] ,
\tag{C.20}
\end{equation*}

对费米子 $n_i, m_j \leq 1$。

最后，用 (C.17) 容易证明：Fock 空间中任何算符的迹为

\begin{equation*}
\operatorname{Tr} \mathcal {A} = \int d ^ {2} \mathbf {z} \, e ^ {- \mathbf {z} ^ {*} \mathbf {z}} \left\langle \mathbf {z} \left| \mathcal {O} \right| \eta \mathbf {z} \right\rangle .
\tag{C.21}
\end{equation*}

注意符号因子 $\eta$——它来自证明过程中 Grassmann 变量次序的调换。

\begin{exercises}

\item 证明正文关于复积分与 Bose 相干态（$\eta=1$）的恒等式：(C.3)、(C.4)、(C.16)、(C.17) 与 (C.19)。

\item 证明关于 Grassmann 积分与 Fermi 相干态（$\eta = -1$）的恒等式：(C.10)、(C.11)、(C.16)、(C.17) 与 (C.19)。

\item 用单位分解 (C.17) 证明：若两个算符在任意两个玻色或费米相干态之间的矩阵元相同，则它们恒等。

\item 指数的正规排序。给定单态的指数双线性算符（Bose 或 Fermi 统计均可）：

\begin{equation*}
B = \exp \left[ \lambda a ^ {\dagger} a \right].
\tag{C.22}
\end{equation*}

证明 Bose 与 Fermi 两种情形都有

\begin{equation*}
\left\langle z \left| B \right| z ' \right\rangle = \exp \left[ e ^ {\lambda} z ^ {*} z ' \right].
\tag{C.23}
\end{equation*}

\item 用上一题与 (C.19) 证明 $B$ 可写成

\begin{equation*}
B = : \exp \left[ \left( e ^ {\lambda} - 1 \right) a ^ {\dagger} a \right] : ,
\tag{C.24}
\end{equation*}

其中 $: O :$ 表示把 $O$ 中算符的每个乘积正规排序。

\item 把 (C.24) 推广到多粒子，证明对厄米矩阵 $\Lambda$：

\begin{align*}
B & = \exp \left[ \mathbf {a} ^ {\dagger} \Lambda \mathbf {a} \right]\\
& = : \exp \left[ \mathbf {a} ^ {\dagger} \left( e ^ {\Lambda} - I \right) \mathbf {a} \right] : .
\tag{C.25}
\end{align*}

（提示：用 $\Lambda$ 的本征基底。）

\end{exercises}

\begin{chapbib}
\item 本附录内容的推荐参考文献：
  J. W. Negele and H. Orland, \textit{Quantum Many Particle Systems} (Addison-Wesley, 1988), 第一章；
  L. S. Schulman, \textit{Techniques and Applications of Path Integration} (Wiley, 1981), 第二十七章。
\end{chapbib}
